\section{Reglas de negocio}

\begin{enumerate}[label*=RN\arabic*]
	
	\item El estado de una celda en el tiempo $t + 1$ depende únicamente de su estado actual y de los estados de sus vecinos.
	
	\item Para iniciar la simulación es necesario que el conjunto de estados, la topología de la vecindad, la condición de frontera y la función de transición hayan sido definidas. 
	
	\item Cuando la simulación está en curso, no se puede cambiar el conjunto de estados, la topología de la vecindad ni la función de transición.
	
	\item El usuario puede manipular el estado de una célula y la condición de frontera en tiempo de ejecución.
	
	\item Los datos importados deben cumplir con el formato definido por el sistema para garantizar su correcta interpretación.
	
	\item Los datos exportados deben conservar el mismo formato que los datos importados.
	
	\item El usuario puede pausar, detener, reanudar o retroceder la simulación sin pérdida de información. 
	
	\item Solo se permite una simulación activa por sesión.
	
\end{enumerate}

\section{Requerimientos}

\subsection{Requerimientos funcionales}

Los requerimientos funcionales describen las capacidades que el sistema debe ofrecer al usuario para cumplir con los objetivos establecidos.

\begin{enumerate}[label=RF\arabic*.]
	
	\item El sistema debe mostrar una cuadrícula regular y finita que represente el espacio celular.
	
	\item El sistema debe permitir al usuario interactuar con el espacio celular, asignando o modificando los estados de las celdas.
	
	\item El sistema debe permitir al usuario definir el conjunto de estados, la topología de la vecindad local y las condiciones de frontera del espacio celular.
	
	\item El sistema debe permitir la evolución del espacio celular mediante un único paso o de forma continua, ajustando la velocidad de actualización según las capacidades computacionales del sistema.
	
	\item El sistema debe permitir la importación y exportación de configuraciones y evoluciones del espacio celular.
	
	\item El sistema debe permitir al usuario establecer la función de transición de manera manual.
	
	\item El sistema debe permitir configurar la red neuronal para imitar el comportamiento de una evolución señalada.
	
	\item El sistema debe entrenar una red neuronal dados los parámetros del usuario y la evolución que desea replicar.
	
\end{enumerate}

\subsection{Requerimientos no funcionales}

Los requerimientos no funcionales especifican las condiciones de calidad, desempeño y compatibilidad bajo las cuales el sistema debe operar.

\begin{enumerate}[label=RNF\arabic*.]
	
	\item El sistema debe poder ejecutarse en Fedora 42 sin requerir modificaciones en el código fuente.
	
	\item El sistema debe ser capaz de procesar al menos 10,000 celdas o registros por segundo utilizando una GPU NVIDIA RTX 5080.
	
	\item El sistema debe garantizar disponibilidad continua durante la ejecución, sin fallos o interrupciones críticas.
	
	\item La interfaz del sistema debe tener mensajes claros y funcionalidades bien definidas.
	
\end{enumerate}

\subsection{Requerimientos de sistema}

\subsubsection{Recomendado}

Los siguientes requerimientos describen el entorno principal y recomendado de ejecución, en el cual el sistema garantiza un rendimiento óptimo y estabilidad durante los procesos de simulación:

\begin{itemize}
	
	\item \textbf{Procesador:} Intel Core i9-14900K × 32 núcleos.
	
	\item \textbf{GPU:} NVIDIA GeForce RTX 5080.
	
	\item \textbf{Sistema operativo:} Fedora Linux 42 (Workstation Edition).
	
	\item \textbf{Entorno de escritorio:} GNOME 48.
	
	\item \textbf{Sistema de ventaneo:} Wayland.
	
\end{itemize}

\subsubsection{Mínimo}

También se establecen los requisitos mínimos necesarios para usar el simulador, teniendo en cuenta que el rendimiento y la velocidad de procesamiento dependerán de las capacidades del hardware disponible y que el entrenamiento podría no ser posible.

\begin{itemize}
	
	\item \textbf{Procesador:} Intel Core i7-1185G7 × 4 núcleos.
	
	\item \textbf{GPU:} Intel Corporation Raptor Lake-S GT1 [UHD Graphics 770].
	
	\item \textbf{Sistema operativo:} Fedora Linux 42 (Workstation Edition).
	
	\item \textbf{Entorno de escritorio:} GNOME 48.
	
	\item \textbf{Sistema de ventaneo:} Wayland.
	
\end{itemize}